
\textbf{Théorème 1.13.}
\emph{Soient \(S\) un schéma noethérien et \(X\) un \(S\)-schéma projectif et muni d'un faisceau ample \({O}_X(1)\) dont les induits \({O}_{X_s}(1)\), avec \(s\in S\), sont engendrés par leur \(H^0({O}_{X_s}(1))\); soit \(\mathcal F\) une famille de classes de faisceaux cohérents sur les fibres de \(X/S\). Alors les conditions suivantes sont équivalentes} :
\begin{enumerate}[label=(\roman*)]
\item \emph{\(\mathcal F\) est limitée. De plus, si tous les \(F_K\in\mathcal F\) sont localement libres de rang \(p\), alors on peut supposer que \(\mathcal F\) est limitée par un faisceau \(F\) sur un \(X_T\), avec \(F\) localement libre de rang \(p\).}
\item \emph{L'ensemble des polynômes de Hilbert \(\chi(F_K(n))\) des \(F_K\in\mathcal F\) est fini, et il existe une suite d'entiers \((b)\) telle que \(\mathcal F\) soit une \((b)\)-famille.}
\item \emph{L'ensemble des polynômes de Hilbert \(\chi(F_K(n))\) des \(F_K\in\mathcal F\) est fini, et il existe un entier \(m\) tel que tous les \(F_K\in\mathcal F\) soient \(m\)-réguliers.}
\item \emph{L'ensemble des polynômes de Hilbert \(\chi(F_K(n))\) des \(F_K\in\mathcal F\) est fini, et \(\mathcal F\) est contenue dans la famille des classes de quotients de faisceaux de la forme \(E_K\), où \(E\) est un faisceau cohérent sur un \(X_T\). De plus, on peut supposer \(T=S\) et \(E\) de la forme \({O}_X(-m)^{\oplus M}\).}
\item \emph{\(\mathcal F\) est contenue dans la famille des classes de conoyaux d'homomorphismes de la forme \(E'_K\to E_K\) où \(E'\) et \(E\) sont des faisceaux cohérents sur un \(X_T\). De plus, on peut supposer \(T=S\) et \(E'\) et \(E\) de la forme \({O}_X(-m')^{\oplus M'}\) et \({O}_X(-m)^{\oplus M}\).}
\end{enumerate}
En effet, supposons que \(\mathcal F\) soit limitée par un faisceau \(F\) sur un \(X_T\); quitte à couper \(T\) en morceaux, en utilisant le théorème de platitude générique (EGA IV 6.9.1), on peut supposer \(F\) plat sur \(T\). Alors,
